%-------------------------
% Statement of Purpose - Amazon ML Summer School 2026
% Ojas Srivastava | Prompt-aligned | Max 500 words
%-------------------------
\documentclass[a4paper,11pt]{article}
\usepackage[empty]{fullpage}
\usepackage{geometry}
\usepackage[hidelinks]{hyperref}
\usepackage{parskip}
\usepackage[T1]{fontenc}

\geometry{left=2.5cm, top=2cm, right=2.5cm, bottom=2cm}
\pagestyle{empty}
\setlength{\parindent}{0pt}
\setlength{\parskip}{10pt}

\begin{document}

\begin{center}
{\LARGE \textbf{Statement of Purpose}} \\[2mm]
{\large Amazon ML Summer School 2026} \\[4mm]
\textbf{Ojas Srivastava} \\
B.Tech Artificial Intelligence, SVNIT Surat \textbar{} CGPA 9.20/10 \textbar{} Graduating 2028 \\
\href{mailto:srivastavaojas454@gmail.com}{srivastavaojas454@gmail.com} \textbar{}
\href{https://github.com/Ojas-Srivastava05}{github.com/Ojas-Srivastava05}
\end{center}

\vspace{2mm}

My journey in AI/ML began at SVNIT Surat, where B.Tech Artificial Intelligence was a deliberate choice to build systems that learn from data. Coursework in Machine Learning, Deep Learning, Probability and Statistics, Linear Algebra, and Optimization gave me theory, but curiosity pushed me further: could models I trained survive noisy data, imperfect labels, and deployment constraints?

What I have built reflects that mindset. On LogiFlow (Google Solution Challenge 2026), I trained Gradient Boosting to predict logistics delays, engineered features, evaluated with train--validation splits, and added Pareto-based ranking over time, cost, and delay before deploying with FastAPI and Redis. On AirHelp (PowerMind Hackathon 2026), I explored embeddings, ChromaDB retrieval, RAG-based QA, and graph pathfinding. At IFFCO, I shipped production Node.js services with Docker and CI/CD. Leading RangRiti and mentoring through ACM SVNIT and Nexus SVNIT taught me to ship products and explain technical ideas clearly.

Competitive programming sharpens this further: LeetCode Knight (contest rating 2048, 711 problems, 32 contests) and Codeforces Specialist (1419, 252+ unique problems solved). That discipline in algorithms and implementation supports both ML coding and the selection test ahead.

These projects also exposed gaps I want MLSS to address. I can train and deploy models, but I need stronger foundations in generalization, probabilistic modeling, optimization for learning, and diagnostics beyond accuracy. I have applied RAG and gradient boosting successfully, yet I want deeper intuition for bias--variance tradeoffs, distribution shift, and principled model selection. I am curious how large-scale systems at Amazon handle evaluation and failure modes---areas hackathon timelines cannot fully cover.

MLSS offers structured modules from Amazon Scientists, assessments on core ML, and exposure to industry-scale problems in search and recommendations. I am especially eager to deepen supervised and unsupervised learning, statistical foundations, and optimization---directly aligned with the program and my remaining undergraduate years.

What makes me a strong fit is how I explore: from coursework to projects, projects to deployment, and failures to better evaluation. I share learning as an ACM executive and mentor and show up consistently in contests and code. I am applying because MLSS matches my curiosity, preparation, and goal of becoming an ML engineer who understands both the mathematics and the systems behind dependable models.

If selected, I will engage fully in every module and assessment. Thank you for considering my application.

\end{document}
